% 英文源分章 PDF 第 14—18 页；第 9.9 节及表 9.1—9.4。
\section{数据分析：弗雷明汉心脏研究}

我们将变量重要性估计方法应用于弗雷明汉冠心病队列。Wilson 等（1998）利用离散化的血压和胆固醇测量值，并结合少量其他变量，开发了按性别分别建立的冠心病风险预测算法。这些风险方程来自按性别分层的多变量回归，用于估计血压、\footnote{血压水平依据 JNC-V（Joint National Committee，1993）定义：理想血压（收缩压 $\leq120$ mm Hg 且舒张压 $\leq80$ mm Hg），正常血压（收缩压 120—129 mm Hg 或舒张压 80—84 mm Hg），正常高值血压（收缩压 130—139 mm Hg 或舒张压 85—89 mm Hg），1 期高血压（收缩压 140—159 mm Hg 或舒张压 90—99 mm Hg），以及 2—4 期高血压（收缩压 $\geq160$ 或舒张压 $\geq100$ mm Hg）。“当收缩压与舒张压落在不同类别时，按较高类别归类。”}胆固醇、年龄、吸烟和糖尿病与未来冠心病（CHD）的关联。胆固醇和血压被分为若干类别，从而能够在线性框架中建模非单调关系。研究者有意采用简单的回归模型，以便临床工作者使用简明的风险评分规则。这些结果至今仍被临床医生广泛用于评估患者的冠心病风险。

我们采用相同的变量和分类离散化方式分析类似数据，但使用实现了数据自适应目标参数 CV-TMLE 方法的 \texttt{varImpact} 软件。具体而言，我们针对原论文的表 5 进行比较分析。与该论文一样，所有分析均按性别分层。（所分析的数据中，女性占 57\%，男性占 43\%。）我们分析的是美国国家心肺血液研究所生物标本与数据资源库信息协调中心提供的公开弗雷明汉纵向数据集（Framingham Longitudinal Data）。分析仅限于收集了低密度脂蛋白（LDL）和高密度脂蛋白（HDL）胆固醇测量值的第 3 期，并删除了在第 1 期或第 2 期已发生 CHD 的受试者。我们只评估原论文所用协变量的影响。

在开始前，先指出我们的数据与原研究存在的几个差别。(1) Wilson 等使用的是弗雷明汉原始队列和子代队列的非公开数据。他们的数据收集于 1971—1974 年，而我们的数据主要收集于 1968 年。(2) 该论文的数据集包含原始队列的子代，而我们的数据集不包含。这可能改变数据生成过程。(3) 1998 年的论文没有说明如何处理缺失数据。我们推测，研究者针对每一张表分别删除了含缺失数据的记录。(4) 为保护患者隐私，弗雷明汉纵向数据集经过匿名化，这很可能对我们的分析结果产生一定影响。因此，我们建议将这些结果视为提示性的，而非定论性的。

\subsection{超级学习器候选库}

我们使用 R 3.3.2 版，通过 \texttt{SuperLearner} 软件包（Polley 和 van der Laan，2013；Polley 等，2017），采用以下超级学习器候选库来估计 $Q$ 和 $g$；该候选库的选择也受到 Pirracchio 等（2014）详尽候选库的启发。该库通过优化结局回归的预测准确性而构建，我们建议在估计变量重要性之前先开展这一有益的工作。

\begin{itemize}
\item 逻辑回归。
\item 弹性网：(1) 仅包含主项；(2) 包含主项和二阶交互项。两种情形均通过 \texttt{glmnet} 软件包（Friedman 等，2010；2.0-5 版）设置六种配置，$\alpha\in\{0,0.2,0.4,0.6,0.8,1.0\}$。
\item 使用 \texttt{arm} 软件包（Gelman 等，2010；1.9-3 版）的贝叶斯线性回归。
\item 使用 \texttt{earth} 软件包（Milborrow 等，2014；4.4.7 版）的多元自适应回归样条，采用三种配置：$\mathrm{degree}\in\{1,2,3\}$。
\item 使用 \texttt{ipred} 软件包（Peters 和 Hothorn，2009；0.9-5 版）的装袋法（bagging）。
\item 使用 \texttt{randomForest} 软件包（Liaw 和 Wiener，2002；4.6-12 版）的随机森林，采用四种配置：$\mathrm{mtry}\in\{1,2,4,7\}$。
\item 使用 \texttt{xgboost} 软件包（0.6-4 版）的极端梯度提升（Chen 和 Guestrin，2016），采用 12 种超参数配置：树的数量 $\in\{100,1000\}$，深度 $\in\{1,2,3\}$，学习率 $\in\{0.1,0.001\}$，取三者的笛卡尔积。
\item 结局均值，用于性能基准比较，并用于检查过拟合。
\end{itemize}

除非另有说明，我们都使用 \texttt{SuperLearner} 软件包提供的默认超参数。为尽可能提高 (1) 统计功效、(2) 随机化假设的可信性，以及 (3) 与原研究的可比性，我们未对调整变量进行任何降维。

\subsection{结果}

下面给出按性别分层的变量重要性结果。第一张表报告所有变量的汇总结果，并在 FDR $p$ 值 $=0.05$ 处添加参考线。第二张表列出每个 CV-TMLE 折的组成参数估计及识别出的水平 $(a_L,a_H)$。我们采用 $V=2$ 的 CV-TMLE。

\noindent\textbf{女性。}从表 9.1 可以看到，HDL 和糖尿病对冠心病的估计影响显著，且在各折之间具有一致性。这意味着，在该数据集的女性中，将 HDL 水平提高到 60+ mg/DL 并预防糖尿病发生，可能降低冠心病风险。这两个结果都与原研究的发现一致。血压（BP）排名最高，但其高、低水平的识别不一致，使其变量重要性的定义更难解释。吸烟状态对 CHD 的影响较小且不显著，这与原研究估计出的较大且显著的效应明显相反。表 9.2 表明，血压的不一致性是由于在一个 CV-TMLE 折中，$a_L$ 被识别为血压“高”水平，而在另一折中，被识别为“正常—理想”水平。HDL 的水平也符合预期：估计高 HDL 的 CHD 发生率最低，低 HDL 的 CHD 发生率最高。相比之下，LDL 在各 CV-TMLE 折间识别出的 $a_L$ 和 $a_H$ 水平确实具有一致性。

\begin{table}[htbp]
\centering
\caption{女性变量重要性结果：合并估计}
\label{tab:ch09-female-combined}
\begingroup
\footnotesize
\setlength{\tabcolsep}{3pt}
\begin{adjustbox}{max width=\textwidth}
\begin{tabular}{@{}rllrcrrc@{}}
\toprule
排名 & 变量 & 类型 & 估计值 & 95\% 置信区间 & $p$ 值 & 调整后 $p$ 值 & 一致性 \\
\midrule
1 & BP & 因子 & 0.1119 & $(0.0388\text{–}0.185)$ & 0.0014 & 0.0081 & 否 \\
2 & HDL & 因子 & 0.1102 & $(0.0257\text{–}0.195)$ & 0.0053 & 0.0125 & 是 \\
3 & 糖尿病 & 有序 & 0.1373 & $(0.0296\text{–}0.245)$ & 0.0062 & 0.0125 & 是 \\
4 & 年龄 & 有序 & 0.0449 & $(-0.00552\text{–}0.0954)$ & 0.0405 & 0.0607 & 是 \\
5 & 吸烟 & 有序 & 0.0122 & $(-0.0362\text{–}0.0605)$ & 0.3108 & 0.3730 & 是 \\
6 & LDL & 因子 & $-0.0193$ & $(-0.0771\text{–}0.0385)$ & 0.7433 & 0.7433 & 否 \\
\bottomrule
\end{tabular}
\end{adjustbox}
\endgroup
\end{table}

\begin{table}[htbp]
\centering
\caption{按估计样本给出的女性变量重要性结果}
\label{tab:ch09-female-folds}
\begingroup
\footnotesize
\setlength{\tabcolsep}{2pt}
\begin{adjustbox}{max width=\textwidth}
\begin{tabular}{@{}lrrllllc@{}}
\toprule
变量 & 估计值 $v1$ & 估计值 $v2$ & 低水平 $v1$ & 高水平 $v1$ & 低水平 $v2$ & 高水平 $v2$ & 一致性 \\
\midrule
BP & 0.1780 & 0.0458 & 高 & 2—4 期 & 正常—理想 & 2—4 期 & 否 \\
HDL & 0.1534 & 0.0670 & $[60,999)$ & $[0,35)$ & $[60,999)$ & $[0,35)$ & 是 \\
糖尿病 & 0.0910 & 0.1836 & $(0.9999999,1]$ & $(1,2]$ & $(0.9999999,1]$ & $(1,2]$ & 是 \\
年龄 & 0.0029 & 0.0870 & $(0.999999,1]$ & $(1,10]$ & $(1,5]$ & $(5,10]$ & 是 \\
吸烟 & 0.0126 & 0.0117 & $(0.9999999,1]$ & $(1,2]$ & $(0.9999999,1]$ & $(1,2]$ & 是 \\
LDL & $-0.0452$ & 0.0067 & $[130,160)$ & $[0,130)$ & $[0,130)$ & $[160,999)$ & 否 \\
\bottomrule
\end{tabular}
\end{adjustbox}
\endgroup
\end{table}

\noindent\textbf{男性。}在男性结果（表 9.3）中，糖尿病和 HDL 再次排名靠前，点估计也与女性相近。这一次，年龄和 LDL 表现出一致、较大且具有统计显著性的效应。吸烟状态的效应仍然较小，不过这次达到边缘统计显著性。对于血压，我们未找到一组一致的高、低水平。表 9.4 中各折的结果表明，男性 HDL 所选择的 $a_L$ 和 $a_H$ 与女性相同。LDL 水平也符合预期：与高 LDL 水平相比，低 LDL 的 CHD 预期风险较低。同样，年龄的 $a_L$ 如预期那样对应最低年龄区间。

\begin{table}[htbp]
\centering
\caption{男性变量重要性结果：合并估计}
\label{tab:ch09-male-combined}
\begingroup
\footnotesize
\setlength{\tabcolsep}{3pt}
\begin{adjustbox}{max width=\textwidth}
\begin{tabular}{@{}rllrcrrc@{}}
\toprule
排名 & 变量 & 类型 & 估计值 & 95\% 置信区间 & $p$ 值 & 调整后 $p$ 值 & 一致性 \\
\midrule
1 & 年龄 & 有序 & 0.1609 & $(0.102\text{–}0.219)$ & 0.0000 & 0.0000 & 是 \\
2 & LDL & 因子 & 0.1693 & $(0.106\text{–}0.232)$ & 0.0000 & 0.0000 & 是 \\
3 & HDL & 因子 & 0.1623 & $(0.09\text{–}0.235)$ & 0.0000 & 0.0000 & 是 \\
4 & 糖尿病 & 有序 & 0.1552 & $(0.0739\text{–}0.236)$ & 0.0001 & 0.0001 & 是 \\
5 & BP & 因子 & 0.0982 & $(0.0384\text{–}0.158)$ & 0.0006 & 0.0008 & 否 \\
6 & 吸烟 & 有序 & 0.0356 & $(-0.0159\text{–}0.0871)$ & 0.0879 & 0.0879 & 是 \\
\bottomrule
\end{tabular}
\end{adjustbox}
\endgroup
\end{table}

\begin{table}[htbp]
\centering
\caption{按估计样本给出的男性变量重要性结果}
\label{tab:ch09-male-folds}
\begingroup
\footnotesize
\setlength{\tabcolsep}{2pt}
\begin{adjustbox}{max width=\textwidth}
\begin{tabular}{@{}lrrllllc@{}}
\toprule
变量 & 估计值 $v1$ & 估计值 $v2$ & 低水平 $v1$ & 高水平 $v1$ & 低水平 $v2$ & 高水平 $v2$ & 一致性 \\
\midrule
年龄 & 0.1339 & 0.1878 & $(0.999999,1]$ & $(1,10]$ & $(0.999999,1]$ & $(1,10]$ & 是 \\
LDL & 0.1139 & 0.2246 & $[0,130)$ & $[160,999)$ & $[0,130)$ & $[160,999)$ & 是 \\
HDL & 0.1542 & 0.1705 & $[60,999)$ & $[0,35)$ & $[60,999)$ & $[0,35)$ & 是 \\
糖尿病 & 0.1897 & 0.1207 & $(0.9999999,1]$ & $(1,2]$ & $(0.9999999,1]$ & $(1,2]$ & 是 \\
BP & 0.1295 & 0.0670 & 正常—理想 & 1 期 & 正常—理想 & 2—4 期 & 否 \\
吸烟 & 0.0333 & 0.0378 & $(0.9999999,1]$ & $(1,2]$ & $(0.9999999,1]$ & $(1,2]$ & 是 \\
\bottomrule
\end{tabular}
\end{adjustbox}
\endgroup
\end{table}
\FloatBarrier
